\subsection{Sprint 2}

Para esta sprint estava previsto que a spike fosse finalizada e que houvesse
respostas definitivas sobre o que seria feito. Além disso, o front-end deveria estar
praticamente pronto, apenas com dados mockados e aguardando as respectivas
rotas do back-end. Por fim, o back-end deveria ter pelo menos um endpoint
rodando, conforme as arquiteturas definidas, com testes implementados. Da minha
parte, fiquei com o back-end, portanto estava previsto finalizar o endpoint de
estados, junto com os testes, e apresentá-lo mediante a devida arquitetura.

O grupo como um todo conseguiu finalizar corretamente as atividades previstas,
entregando uma apresentação definitiva das tecnologias a serem utilizadas para o
cliente, além de entregar o front-end com alguns detalhes, mas em geral totalmente
implementado. Por fim, o back-end também conseguiu evoluir bem: já tínhamos um
endpoint com router, service e repository configurados e testes de integração
feitos.

Diferente da minha primeira AGES, a tecnologia utilizada, Python, não era de tanto
conhecimento da equipe e, portanto, passamos por alguns apertos. Práticas como o
mapeamento correto do banco por meio do SQLAlchemy e a criação de funções
assíncronas no back-end exigiram mais estudo. Além disso, a definição de uma
arquitetura foi um pouco mais complicada do que o normal, considerando as
diferentes nomenclaturas que o Python adota e suas diferentes práticas.

Apesar de terem ocorrido problemas, eles foram resolvidos e se tornaram
aprendizados. Foi possível perceber que utilizar uma linguagem em que ninguém do
time tinha grande conhecimento era um desafio importante e exigia reflexão sobre
seu custo. Ao mesmo tempo, cada bug resolvido, erro percorrido e problema
enfrentado desde a configuração até o código virou experiência que, se bem
absorvida, pode facilitar trabalhos futuros.
